% Generated draft: requires mathematical and source review.
\documentclass[11pt]{article}
\usepackage[margin=25mm]{geometry}
\usepackage{fontspec}
\setmainfont{Latin Modern Roman}
\usepackage{amsmath,amssymb,mathtools,unicode-math}
\setmathfont{Latin Modern Math}
\usepackage{enumitem,needspace,xurl,hyperref}
\setlength{\parindent}{0pt}
\setlength{\parskip}{5pt}
\setlength{\emergencystretch}{3em}
\newcommand{\N}{\mathbb{N}}
\newcommand{\Z}{\mathbb{Z}}
\newcommand{\Q}{\mathbb{Q}}
\newcommand{\R}{\mathbb{R}}
\newcommand{\C}{\mathbb{C}}
\newcommand{\F}{\mathbb{F}}
\newcommand{\A}{\mathbb{A}}
\newcommand{\PP}{\mathbb{P}}
\newcommand{\E}{\mathbb{E}}
\newcommand{\eps}{\varepsilon}
\newcommand{\dd}{\,\mathrm{d}}
\begin{document}
\section*{UnsolvedMath ID: 152}
\textbf{Source problem code:} catalogue entry
\subsection{Definitions and mathematical statement}
A group \(G\), with identity element \(1_G\), is \emph{residually finite} if
\[
\forall g\in G\setminus\{1_G\}\ \exists\text{ a finite group }F\ \exists\text{ a group homomorphism }\varphi:G\to F
\quad \varphi(g)\ne1_F.
\]
Equivalently, each nonidentity element survives in some finite quotient. Under the residual-finiteness interpretation supplied by the title and background, the question ``Is every group well-approximated by finite groups?'' is the universal assertion
\[
\forall\text{ groups }G\quad G\text{ is residually finite}.
\]
No finite-generation or finite-presentation hypothesis occurs in this assertion.

The supplied literature triage dated 2026-08-17 classifies this universal assertion as false under that interpretation, identifying nontrivial infinite simple groups as a counterexample class and citing Higman (1951). The registry's separate status field is open; this is retained as source metadata rather than a mathematical conclusion. The background speculates about finitely generated or finitely presented variants, but does not actually formulate one; in particular it cannot silently replace the universal domain above. The supplied Higman reference is described by the triage as providing a finitely generated infinite simple group, so finite generation by itself is not an asserted escape from that counterexample class.

If ``well-approximated'' means an approximation notion other than residual finiteness, the source does not define it. Such a reading would be a different mathematical question and is not resolved by the conditional formulation given here.
\subsection{Short English statement}
Under the source's residual-finiteness interpretation, every nonidentity element would have to remain nonidentity in some finite image, for every group. The supplied 2026-08-17 triage reports this universal assertion false using infinite simple groups. An alternative meaning of finite approximation is not specified.
\subsection{Sources}
\begin{itemize}
\item Local UnsolvedMath registry, permanent ID 152, source code catalogue entry, ``Residually Finite Groups,'' including the supplied literature triage dated 2026-08-17. Membership in Ben Green's problem set is registry metadata; no direct Green problem reference is supplied for this wording.
\item G. Higman, ``A Finitely Generated Infinite Simple Group,'' \emph{Journal of the London Mathematical Society} 26 (1951), 61--64, supplied primary reference: \url{https://londmathsoc.onlinelibrary.wiley.com/doi/pdf/10.1112/jlms/s1-26.1.61}. The supplied triage cites this for the counterexample class under the residual-finiteness reading.
\end{itemize}
\end{document}
